
\section{代数预备}

设 $A$ 是一个\trRing，而 $J_1\supset J_2\supset\cdots$ 是一个递减的\trIdeal{}序列。极限 $\lim A/J_n$ 常被记作 $\hat{A}$，并称为 $A$ 的\trCompletion（针对这一列\trIdeal）；这列\trIdeal{}则被隐去不言。\trCompletion{}最重要的情形是 $J_n=I^n$，其中 $I$ 是 $A$ 的一个\trIdeal。极限 $\lim A/I^n$ 称为\textbf{$A$ 的 $I$-进\trCompletion}（$I$-adic completion of $A$），或\textbf{$A$ 沿 $I$ 的\trCompletion}或\textbf{在 $I$ 处的\trCompletion}。

\textbf{28.1.1. \trhExample。}\quad 我们把 $k[[x_1,\dots,x_n]]$ 定义为 $k[x_1,\dots,x_n]$ 在\trMaximalIdeal{} $(x_1,\dots,x_n)$ 处的\trCompletion。这是 $k$ 上 $n$ 个变量的\trFormalPowerSeries{}环（\mbox{例~1.3.3}）。

\trPadics{}（\mbox{例~1.3.4}）是另一个例子。

\textbf{28.1.A.（练习）}\quad 设 $J'_1\supset J'_2\supset\cdots$ 是一个与 $J_n$ \textbf{共尾}（cofinal）的递减\trIdeal{}序列。换言之，对每个 $J_n$，存在某个 $J'_N$ 使得 $J_n\supset J'_N$；而对每个 $J'_n$，存在某个 $J_N$ 使得 $J'_n\supset J_N$。证明：存在\trCanonical{}\trIsomorphism
\[
\lim A/J'_n \xleftrightarrow{\ \sim\ } \lim A/J_n.
\]
于是，真正要紧的与其说是\trIdeal{}序列的具体样貌，不如说是其「共尾」\trEquivalenceClass。

\textbf{28.1.2. 预备注记。}\quad 我们有一个显然的\trMorphism{} $A\to\hat{A}$。

若在 $A/J_n$ 上赋予\trDiscreteTopology，则 $\hat{A}$ 自然而然地具有\textbf{拓扑环}（topological ring）的结构（所谓拓扑环，就是同时是一个\trTopologicalSpace{}的\trRing，其中一切环运算都是连续的；等价地说，即\trTopologicalSpace{}\trCategory{}中的一个\trRing{}对象）。于是我们还可以谈论拓扑环 $A'$ 上的\textbf{拓扑模}（topological module）$M$——即底\trRing{} $A'$ 上的一个\trModule，带有拓扑，使得 $A'$ 在 $M$ 上的作用是连续的。这是一种有用的观点，但我们不会用到它。

\textbf{28.1.B.（练习）}\quad 设 $\mathfrak{m}$ 是环 $A$ 的一个\trMaximalIdeal。证明：$A$ 在 $\mathfrak{m}$ 处的\trCompletion{}\trCanonical{}地\trIsomorphic{}于 $A_{\mathfrak{m}}$ 在 $\mathfrak{m}$ 处的\trCompletion。

若 $(A,\mathfrak{m})$ 是\trNoether{}\trLocalRing，则自然映射 $A\to\hat{A}$ 是\trInjectiveMap：\trKernel{}中的任何元素都必定落在 $\bigcap\mathfrak{m}^i$ 中，而由\trKrull{}交定理（\mbox{练习~13.9.A}(b)），后者为 $0$。于是「完备化不丢失任何信息」，正如（局部地看）解析函数由其幂级数展开所决定。

（在\trIdeal{} $I$ 处取\trCompletion{}的情形，若 $A\to\hat{A}$ 是\trInjectiveMap，则称 $A$ 是 \textbf{$I$-进分离}的（$I$-adically separated）；若 $A\to\hat{A}$ 是\trIsomorphism，则称 $A$ \textbf{对 $I$ 完备}（complete with respect to $I$）——尽管我们用不到这些措辞。例如，\trKrull{}交定理蕴含：若 $I$ 是\trNoether{}\trIntegralDomain{}或\trNoether{}\trLocalRing{}的一个\trProperIdeal，则 $A$ 是 $I$-进分离的。）

\textbf{28.1.C.（练习）}\quad 设 $(A,\mathfrak{m})$ 是一个包含其\trResidueField{} $k$ 的\trNoether{}\trLocalRing（也就是说，它是一个 $k$-\trAlgebra）。设 $x_1,\dots,x_n$ 是 $A$ 的若干元素，其在 $\mathfrak{m}/\mathfrak{m}^2$ 中的像构成一组基。证明：$k$-代数映射
\begin{equation*}
\tag{28.1.2.1}\label{eq:formal-series-to-completion}
k[[t_1,\dots,t_n]]\longrightarrow\hat{A}
\end{equation*}
（由 $t_i\mapsto x_i$ 定义）是满射。（请先解释为什么\emph{存在}这样的映射！）照例，对\trLocalRing{}而言，\trCompletion{}默认取在\trMaximalIdeal{}处。

\mbox{练习~28.1.C} 是\emph{\trCohen{}结构定理}（Cohen Structure Theorem）的特例。（关于这一话题的更多内容，见 [E, \S7.4]。）

\textbf{28.1.D.（练习）}\quad 设 $X$ 是 $k$ 上的一个\trLocally\trNoether{}\trScheme，而 $p\in X$ 是一个\trRationalPoint{}（$k$-\trValuedPoint）。设 $p\in X$ 是\trRegular{}的，\trCodimension{}为 $d$（即 $\dim\OO_{X,p}=d$）。请描述一个作为拓扑环的\trIsomorphism{} $\hat{\OO}_{X,p}\xleftrightarrow{\ \sim\ }k[[x_1,\dots,x_d]]$。（提示：如\mbox{练习~28.1.C} 中那样，选取 $\mathfrak{m}$ 的 $d$ 个元素，使其在 $\mathfrak{m}/\mathfrak{m}^2$ 中的像构成一组基；它们将是你的 $x_1,\dots,x_d$。证明：映射 \eqref{eq:formal-series-to-completion} 没有核。利用\mbox{定理~22.3.8} 把 $\mathfrak{m}^n/\mathfrak{m}^{n+1}$ 等同于 $\Sym^n(\mathfrak{m}/\mathfrak{m}^2)$ 或许有帮助。）

逆命题同样成立：若 $(A,\mathfrak{m},k)$ 是一个作为 $k$-代数的\trNoether{}\trLocalRing，且 $\hat{A}\cong k[[x_1,\dots,x_d]]$，则 $A$ 是一个 $d$ 维的\trRegular\trLocalRing；关键一步见 [AtM, Prop.~11.24]。

\textbf{28.1.3. \trhRemark。}\quad 设 $p$ 是 $d$ 维 $k$-\trVariety{}上的一个光滑的有理（$k$-值）点，并设 $f\in\OO_{X,p}$（$f$ 是一个「局部函数」）。借助\mbox{练习~28.1.D} 的\trIsomorphism，我们可以把 $f$ 解读为 $k[[x_1,\dots,x_d]]$ 中的一个元素。这应当被解读为「$f$ 在 $p$ 处的幂级数展开」，展开于「局部坐标 $x_1,\dots,x_d$」之下。

\textbf{28.1.4. 一些更多的几何动机。}\quad \mbox{练习~28.1.D} 或许能为\trCompletion{}给出一点动机：「在\trCompletion{}中，\trRegular\trScheme{}看起来像仿射空间」。这句话常以富有暗示性的语言陈述为「形式局部地，正则概形同构于仿射空间」；\S28.3.1 将把这一点说得略微更精确些。

我们现在再给出一点\trCompletion{}的几何动机，后面不会用到它。回想 \S13.6.1：平展\trMorphism{}被设计得看起来像微分几何中的「局部同构」。但\mbox{练习~13.6.F} 表明，这个比喻在一个重要方面败得很惨。更精确地说：令 $Y=\Spec k(t)$，$X=\Spec k[u,1/u]$，令 $p\in X$ 由 $u=1$ 给出，而 $q\in Y$ 由 $t=1$ 给出。设 $\operatorname{char}k\ne 2$。由 $t\mapsto u^2$ 诱导的\trMorphism{} $\pi\colon X\to Y$ 是平展的，且 $\pi(p)=q$。但 $p$ 不存在这样的开邻域，使得 $\pi$ 把它\trIsomorphic{}地映到 $q$ 的某个开邻域上。然而，下面的练习将表明：$\pi$ 诱导\trCompletion{}之间的\trIsomorphism（即「形式邻域」之间的\trIsomorphism）。

\textbf{28.1.E.（练习）}\quad 沿用上一段的记号，证明：$\pi$ 诱导\trCompletion{}之间的\trIsomorphism{} $\hat{\OO}_{Y,q}\to\hat{\OO}_{X,p}$。

\textbf{28.1.F.（练习）}\quad 设 $Y=\Spec A$，$X=\Spec A[t]/(f(t))$，而 $\pi\colon X\to Y$ 是由 $A\to A[t]/(f(t))$ 诱导的\trMorphism。设 $p\in X$ 是\trClosedPoint，$q\in Y$ 是\trClosedPoint，$\pi(p)=q$，且 $\pi$ 在 $p$ 处平展并诱导\trResidueField{}之间的\trIsomorphism。证明：$\pi$ 诱导\trCompletion{}之间的\trIsomorphism{} $\hat{\OO}_{Y,q}\to\hat{\OO}_{X,p}$。（如果你还不熟悉\textbf{\trHensel{}引理}（Hensel's Lemma），你或许会在解这道题的过程中重新发现它的核心思想。）

稍加细心，你便可以更一般地证明：若 $\pi\colon X\to Y$ 是平展\trMorphism，$\pi(p)=q$，且 $\pi$ 在 $p$ 处诱导\trResidueField{}之间的\trIsomorphism，则 $\pi$ 诱导\trCompletion{}之间的\trIsomorphism{} $\hat{\OO}_{Y,q}\to\hat{\OO}_{X,p}$。（你甚至可能想琢磨一下，如何去掉\trResidueField{}之间\trIsomorphism{}这一假设。）

这些想法与 \S24.8.13 中讨论的「形式\trEtale」定义相去不远。

你可以把这些结果解读为这样的陈述：「隐函数定理在形式局部上管用」，尽管它在扎里斯基局部上并不管用。（平展\trTopology{}介于这两者之间；它的部分设计意图，正是让这种意义下的隐函数定理总是成立——这纯粹是约定使然。但此处不宜讨论平展\trTopology。）

下面还有两个有趣的问题，它们暗地里讲的都是\trHensel{}引理。

\textbf{28.1.G.（练习）}\quad 证明：$x^2+y^2=-2$ 在模 $7^{100}$ 之下有解。

\textbf{28.1.H.（练习）}\quad 设 $\mathbb{A}^2_{\mathbb{Q}}$ 中有一条平面\trCurve{} $f(x,y)=0$，而 $(0,a)$ 是它的一个光滑点，其\trTangentLine{}非竖直。证明：恰好存在一个幂级数解 $y=p(x)$，使得 $f(x,p(x))=0$（作为幂级数），并且 $p(0)=a$。

\textbf{28.1.5. \trCompletion{}与\trExactness。}\quad 我们用一个有趣而有用的陈述来结束本节。

\textbf{28.1.6. \trhTheorem。}\ ——\ \emph{设 $A$ 是\trNoether{}环，而 $I\subset A$ 是一个\trIdeal。对任意 $A$-\trModule{} $M$，令 $\hat{M}=\lim M/I^jM$ 为 $M$ 关于 $I$ 的\trCompletion。}

\emph{(a) \trCompletion{} $\hat{A}$（$A$ 关于 $I$ 的\trCompletion）在 $A$ 上\trFlat。}

\emph{(b) 若 $M$ 是有限生成的，则自然映射 $\hat{A}\otimes_A M\to\hat{M}$ 是\trIsomorphism。}

\emph{(c) 若 $0\to M\to N\to P\to 0$ 是有限生成 $A$-\trModule{}的一个\trShortExactSequence，则 $0\to\hat{M}\to\hat{N}\to\hat{P}\to 0$ 正合。（于是，由\mbox{练习~1.5.F}，\trCompletion{}保持有限生成\trModule{}的\trExactSequence。）}

我们将在 \S28.2 中用到 (a)，在 \S28.7 中用到 (b)。

\textbf{28.1.7. \trhRemark。}\quad 在证明\mbox{定理~28.1.6} 之前，我们先作一些注记。(a) 与 (b) 合在一起显然蕴含 (c)，但我们却要用 (c) 来证明 (a) 与 (b)。另外，注意一个微妙的区别（它帮助我记住这个陈述）：若 $0\to M\to N\to P\to 0$ 是 $A$-\trModule{}的一个\trExactSequence，其中的模\emph{未必有限生成}，则
\begin{equation*}
\tag{28.1.7.1}\label{eq:completion-tensor-exact}
0\longrightarrow \hat{A}\otimes_A M\longrightarrow \hat{A}\otimes_A N\longrightarrow \hat{A}\otimes_A P\longrightarrow 0
\end{equation*}
\emph{总是}正合，但
\begin{equation*}
\tag{28.1.7.2}\label{eq:completion-itself-exact}
0\longrightarrow \hat{M}\longrightarrow \hat{N}\longrightarrow \hat{P}\longrightarrow 0
\end{equation*}
\emph{未必}正合——而当它\emph{的确}正合时，往往是因为这些\trModule{}是有限生成的，从而 \eqref{eq:completion-itself-exact} 其实就是 \eqref{eq:completion-tensor-exact}。

当心：\trCompletion{}并不总是正合的。考虑 $k[t]$-\trModule{}的正合列
\[
0\longrightarrow \bigoplus_{n=1}^{\infty}k[t]\xrightarrow{\ \times(t,t^2,t^3,\dots)\ }\bigoplus_{n=1}^{\infty}k[t]\longrightarrow \bigoplus_{n=1}^{\infty}k[t]/(t^n)\longrightarrow 0.
\]
对理想 $((t,t,t,\dots))$ 作\trCompletion{}之后，这个序列在中间不再正合：$(t^2,t^3,t^4,\dots)$ 被映到 $0$，但它并不在前一项的\trCompletion{}的像中。

\extrahard\ \emph{证明。}\quad 关键的一步是证明 (c)，我们通过一连串练习来完成它。设 $0\to M\to N\to P\to 0$ 是有限生成 $A$-\trModule{}的一个\trShortExactSequence。

\textbf{28.1.I.（练习）}\quad 证明：$\hat{N}\to\hat{P}$ 是满射。提示：把 $\hat{P}$ 的一个元素看成序列 $(p_j\in P/I^jP)_{j\ge 0}$，其中 $p_{j+1}$ 的像是 $p_j$；参看\mbox{练习~1.3.C}。对 $j$ 作归纳，构造出原像 $(n_j\in N/I^jN)_{j\ge 0}$。

我们现在希望把 $\ker(\hat{N}\to\hat{P})$ 等同于 $\hat{M}$。

\textbf{28.1.J.（练习）}\quad 证明：对每个 $j\ge 0$，
\begin{equation*}
\tag{28.1.7.3}\label{eq:ses-mod-intersection}
0\longrightarrow M/(M\cap I^jN)\longrightarrow N/I^jN\longrightarrow P/I^jP\longrightarrow 0
\end{equation*}
是正合的。（一个可能的提示：证明 $0\to M\cap I^jN\to M\to N/I^jN\to P/I^jP\to 0$ 正合。）

这些短正合列 \eqref{eq:ses-mod-intersection} 随着 $j$ 的变动构成一个\trInverseSystem。它的极限是\trLeftExact{}的（极限总是如此），但由\mbox{练习~1.5.N}，它还是\trRightExact{}的，因为「左边的\trTransitionMap{}」$M/(M\cap I^{j+1}N)\to M/(M\cap I^jN)$ 显然是满射。于是
\begin{equation*}
\tag{28.1.7.4}\label{eq:limit-ses}
0\longrightarrow \lim M/(M\cap I^jN)\longrightarrow \hat{N}\longrightarrow \hat{P}\longrightarrow 0
\end{equation*}
正合。为了完成 (c) 的证明，只需（由\mbox{练习~28.1.A}）证明：子模序列 $I^jM$ 与序列 $M\cap I^jN$ 共尾，从而 \eqref{eq:limit-ses} 左端的项 $\lim M/(M\cap I^jN)$ 便自然而然地等同于 $\hat{M}$。

\textbf{28.1.K.（练习）}\quad 证明这一点。提示：显然，$I^jM\subset M\cap I^jN$。由\trArtin--\trRees{}引理~13.9.3 的\trCorollary{}~13.9.4，对某个整数 $s$，对一切 $j\ge 0$ 都有 $M\cap I^{j+s}N=I^j(M\cap I^sN)$，而显然 $I^j(M\cap I^sN)\subset I^jM$。

这就完成了\mbox{定理~28.1.6} 的部分 (c) 的证明。

对于部分 (b)，把 $M$ 呈现为
\begin{equation*}
\tag{28.1.7.5}\label{eq:presentation-of-m}
A^{\oplus m}\xrightarrow{\ \alpha\ }A^{\oplus n}\longrightarrow M\longrightarrow 0,
\end{equation*}
其中 $\alpha$ 是一个系数取在 $A$ 中的 $m\times n$ 矩阵。在这一情形下，由部分 (c)，\trCompletion{}是正合的，而且它显然与有限\trDirectSum{}交换，于是
\[
\hat{A}^{\oplus m}\longrightarrow \hat{A}^{\oplus n}\longrightarrow \hat{M}\longrightarrow 0
\]
是正合的。\trTensorProduct{}是\trRightExact{}的，并且与\trDirectSum{}交换（\mbox{练习~1.2.M}），于是
\[
\hat{A}^{\oplus m}\longrightarrow \hat{A}^{\oplus n}\longrightarrow \hat{A}\otimes_A M\longrightarrow 0
\]
也是正合的。注意，这两列\trRightExact{}的序列中从 $\hat{A}^{\oplus m}$ 到 $\hat{A}^{\oplus n}$ 的映射是相同的；它们都是 $\alpha$。于是二者的\trCokernel{}得到等同，(b) 随之得证。

最后，为了证明 (a)，我们需要把平坦性的理想论判别法（\mbox{定理~24.3.1}）稍作推广。回想（\S24.3.2）：它等价于下述事实——$A$-\trModule{} $M$ 是\trFlat{}的，当且仅当对一切\trIdeal{} $I$，自然映射 $I\otimes_A M\to M$ 都是\trInjectiveMap。

\textbf{28.1.L.（练习：平坦性的理想论判别法的加强形式）}\quad 证明：$A$-\trModule{} $M$ 是\trFlat{}的，当且仅当对一切\emph{有限生成}的\trIdeal{} $I$，自然映射 $I\otimes_A M\to M$ 都是\trInjectiveMap。（提示：若对某个非有限生成的\trIdeal{} $J$ 有反例，就用它找出另一个对\emph{确为}有限生成的\trIdeal{} $I$ 的反例。）

依这条判别法，为证 (a)，只需证明：对一切有限生成的\trIdeal{} $I$，乘法映射 $I\otimes_A\hat{A}\to\hat{A}$ 都是\trInjectiveMap。但由部分 (b)，这与证明 $\hat{I}\to\hat{A}$ 是\trInjectiveMap 是同一回事\footnote{译注：原书此处作「this is the same showing that」，当为「this is the same as showing that」之误（脱漏 as），径改。——译者注}；而这由部分 (c) 即得。\hfill$\square$
